\subsection{Requerimientos más importantes, con su justificación}
\label{anx:requerimientos}

Esta sección amplía la Tabla~\ref{tab:req-resumen} con la descripción y la justificación de los requerimientos más relevantes del sistema, organizados en no
funcionales y funcionales. Cada requerimiento se identifica con un código propio de esta
sección, seguido de su nombre y, en el caso de los no funcionales, de la característica de
calidad de la norma ISO/IEC 25010 a la que pertenece \citep{iso_25010_2011}.\par

\paragraph{Requerimientos no funcionales.} Estos requerimientos definen cómo debe comportarse
el sistema. Los más importantes se relacionan con la protección de los datos de las familias
y de los menores, con la seguridad de los pagos y con el tiempo de respuesta, ya que de ellos
depende la confianza del acudiente en la plataforma.\par

\begin{itemize}
  \item \textbf{\textit{RNF01 -- No almacenar datos de pago (Seguridad)}}
    \begin{itemize}
      \item \textit{Descripción:} el sistema no debe guardar números de tarjeta ni datos
            bancarios. La captura del medio de pago debe ocurrir por completo en la pasarela
            de pagos.
      \item \textit{Justificación:} protege la información financiera de las familias y
            reduce el riesgo de que una falla de seguridad en Tukit exponga datos de pago.
    \end{itemize}

  \item \textbf{\textit{RNF02 -- Protección de los datos de los menores (Seguridad)}}
    \begin{itemize}
      \item \textit{Descripción:} el sistema debe guardar solo los datos indispensables del
            estudiante, cifrarlos y eliminar su nombre de la lista antes de enviarla al
            modelo de lectura, que es un servicio externo.
      \item \textit{Justificación:} la Ley 1581 de 2012 y el Decreto 1377 de 2013 dan un
            tratamiento especial a los datos de niños y adolescentes
            \citep{congreso_ley1581_2012,presidencia_decreto1377_2013}.
    \end{itemize}

  \item \textbf{\textit{RNF03 -- Tiempo de generación de los kits (Eficiencia de
        desempeño)}}
    \begin{itemize}
      \item \textit{Descripción:} una vez normalizada la lista, el sistema debe generar los
            kits en 20 segundos o menos en el 95\,\% de los casos, y el proceso completo,
            desde la carga de la lista hasta la creación de los kits, no debe superar los
            cinco minutos.
      \item \textit{Justificación:} el acudiente espera este resultado frente a la pantalla,
            y una espera excesiva lleva a abandonar el proceso.
    \end{itemize}

  \item \textbf{\textit{RNF04 -- Acceso desde el celular sin instalación (Usabilidad)}}
    \begin{itemize}
      \item \textit{Descripción:} la aplicación del acudiente debe funcionar desde el
            navegador del celular como una aplicación web progresiva, sin necesidad de
            instalarla.
      \item \textit{Justificación:} la mayoría de los acudientes compra desde el celular, y
            evitar la instalación elimina una barrera para el primer uso.
    \end{itemize}



\item \textbf{\textit{RNF05 -- Capacidad en temporada escolar (Eficiencia de
      desempeño)}}
  \begin{itemize}
    \item \textit{Descripción:} el sistema debe atender a 50 usuarios simultáneos y del
          orden de 1.000 pedidos en el día de mayor demanda, sin que los tiempos de
          respuesta se deterioren de forma significativa y con una tasa de error inferior
          al 5\,\%.
    \item \textit{Justificación:} la demanda se concentra en pocas semanas del año. La
          meta equivale a diez veces los cerca de 95 pedidos diarios que se proyectan para
          la temporada del primer año, lo que deja un margen amplio frente a la demanda
          esperada.
  \end{itemize}

\item \textbf{\textit{RNF06 -- Crecimiento sin rediseño (Escalabilidad)}}
  \begin{itemize}
    \item \textit{Descripción:} la arquitectura debe permitir crecer hasta 1.000 usuarios
          simultáneos aumentando el número de copias de cada servicio, sin cambiar su
          diseño ni su modelo de datos.
    \item \textit{Justificación:} el sistema debe poder incorporar nuevos colegios y
          ciudades después del piloto sin reconstruir la plataforma.
  \end{itemize}

\item \textbf{\textit{RNF07 -- Disponibilidad en las horas de mayor uso (Fiabilidad)}}
  \begin{itemize}
    \item \textit{Descripción:} el sistema debe estar disponible al menos el 90\,\% del
          tiempo durante las franjas de mayor demanda, de 12:00 a 14:00 y de 18:00 a 22:00.
          El mantenimiento programado se realiza entre las 02:00 y las 03:00.
    \item \textit{Justificación:} en esas franjas se concentra la actividad de las familias,
          y una interrupción en temporada escolar puede llevar al acudiente a comprar por
          otro medio.
  \end{itemize}

\item \textbf{\textit{RNF08 -- Tolerancia a fallas (Fiabilidad)}}
  \begin{itemize}
    \item \textit{Descripción:} la falla de un servicio no debe afectar a los demás. Si el
          servicio externo de lectura deja de responder, el sistema debe seguir mostrando
          el catálogo y los kits ya generados, y ofrecer al acudiente la opción de ingresar
          su lista manualmente.
    \item \textit{Justificación:} la lectura depende de un proveedor externo que puede
          fallar, y esa falla no debe impedir que las familias consulten o compren lo que
          ya está listo.
  \end{itemize}
\end{itemize}

\paragraph{Requerimientos funcionales.} Estos requerimientos describen lo que el sistema
debe hacer, y se agrupan según el recorrido de una lista, desde su carga hasta la compra.
Se enfocan en los tres actores que usan la plataforma: el acudiente, el colegio y el equipo
de Tukit.\par

\begin{itemize}
  \item \textbf{\textit{RF01 -- Cargar y leer la lista}}
    \begin{itemize}
      \item \textit{Descripción:} el acudiente debe poder seleccionar la lista oficial de su
            colegio y grado, o cargar la suya como fotografía, PDF o texto, y el sistema
            debe convertirla en renglones con su texto y su cantidad.
      \item \textit{Justificación:} cada colegio publica su lista en un formato distinto, y
            sin esta capacidad el acudiente tendría que transcribirla a mano.
    \end{itemize}

  \item \textbf{\textit{RF02 -- Reconocer las distintas formas de nombrar un artículo}}
    \begin{itemize}
      \item \textit{Descripción:} el sistema debe asociar cada renglón con su concepto del
            catálogo, aunque esté escrito con sinónimos o errores de ortografía, y enviar a
            revisión los renglones que no logre resolver con certeza suficiente.
      \item \textit{Justificación:} las listas de diecisiete colegios nombran los útiles de
            1.451 formas que corresponden a 99 conceptos, y sin este mecanismo la mayoría de
            los renglones quedaría sin producto.
    \end{itemize}

  \item \textbf{\textit{RF03 -- Generar los kits}}
    \begin{itemize}
      \item \textit{Descripción:} el sistema debe generar hasta cuatro kits completos
            (Ahorrador, Premium, Fan y A tu Medida), con un máximo de tres opciones de
            producto por renglón.
      \item \textit{Justificación:} ofrecer pocas alternativas bien diferenciadas facilita la
            decisión del acudiente \citep{iyengar_when_2000}.
    \end{itemize}

  \item \textbf{\textit{RF04 -- Personalizar y pagar el kit}}
    \begin{itemize}
      \item \textit{Descripción:} el acudiente debe poder cambiar cualquier producto, elegir
            el kit que prefiera, pagarlo en línea y solicitar el retracto dentro del plazo
            legal.
      \item \textit{Justificación:} la recomendación orienta la decisión sin imponerla, y la
            Ley 1480 de 2011 obliga a garantizar el retracto en las compras a distancia
            \citep{congreso_ley1480_2011}.
    \end{itemize}

  \item \textbf{\textit{RF05 -- Publicar y revisar las listas oficiales}}
    \begin{itemize}
      \item \textit{Descripción:} el colegio debe poder publicar la lista oficial de cada
            grado, y el equipo de Tukit debe revisarla y validar sus renglones antes de que
            llegue a las familias.
      \item \textit{Justificación:} garantiza que las familias compren sobre la versión
            oficial de la lista y no sobre copias desactualizadas.
    \end{itemize}
\end{itemize}
